% A new local action. Every preceding source and numerical artifact is unchanged.
\subsection{Restoring a fourth node before its own even cube test}
\label{sec:lower-fourth-own-cube-restoration}

The previous fourth refinement first tested
$S_3+3t_4<\lambda_3$, then changed the cap only at admitted nodes.
That is a restriction of that particular action, not a prohibition
imposed by the lower parity identity. There is a different legal
action: at a selected fourth frontier prefix, choose its licensed
cap before testing that fourth node's own cube. This note proves the
action and one exact positive rational box. It does not recalculate
the global functional or append a numerical gain to a frozen certificate.

Write $\nu=249/448$, $h=19/125$, $\Lambda=153/256$ and
$p_j=N^{t_j}$, $t_1>\cdots>t_4>0$. Keep the actual third ancestor,
its branch and its frozen cap $\lambda_3$ unchanged. The earlier
history, including $t_1+3t_2<\nu$ (or its buffered version), must
still hold. Select finitely many fourth boxes that satisfy
\[
 S_3+3t_4\ge\lambda_3
\]
throughout, so these nodes are old first exits. Choose a fixed new
level $\lambda_4$ on each selected box such that, with strict margins,
\begin{equation}\label{eq:own-cube-new-tests}
 \lambda_4>\lambda_3,\qquad S_3+3t_4<\lambda_4,
 \qquad \max_{j\le4}(S_{j-1}+2t_j)<\lambda_4,
 \qquad G_I<\lambda_4<v(b_I).
\end{equation}
Here $b_I=\sum_{i\in I}t_i$ and
\[
 G_I=\max_{j\notin I}\left(
 b_I+\sum_{\substack{i<j\\i\notin I}}t_i+2t_j\right),
 \qquad
 v(b)=\min\{(1+b)/2,(32-34b)/39,57/64-3b/2\}.
\]
The inequalities are new whole-box support tests; one cannot claim
them merely because the third ancestor was licensed at its old level.
The new decision depends only on the first four primes and fixed
policy data, not on a future prime or a progression error.

\paragraph{Pointwise lower parity, not monotonicity of signed weights.}
At a restored fourth prefix $P_4$ retain its coefficient $+1$ and
attach the ordinary residual lower support
\[
 \mathcal D^-(R),\qquad R=N^{\lambda_4}/P_4,
 \qquad \text{residual primes }q<p_4.
\]
The own-cube test gives $R>p_4^2$. Consequently the residual first
square $q_1^2<R$ never fails. All later first exclusions have even
residual rank, hence even full rank $6,8,\ldots$. At unselected
fourth frontier nodes retain the old rule. The resulting tree is
prefix closed and all its first exits still have even full rank.

For any finite set of primes $Q$ and $0\le q(p)\le1$, partition
excluded squarefree divisors by their unique first excluded prefix.
Since each such prefix has even rank, the exact identity is
\[
 \sum_{d\mid Q}\mu(d){\bf1}_{d\in\mathcal T}q(d)
 =\prod_{p\mid Q}(1-q(p))
  -\sum_{e\text{ first exit}}q(e)
                \prod_{p< p_{\rm last}(e),\ p\mid Q}(1-q(p)).
\]
Take $q(p)={\bf1}_{p\mid n}$. This proves a pointwise lower weight.
It does not say that the newly added signed subtree is nonnegative
at every integer, or that finite signed weights increase pointwise.
The exact raw-density difference is instead
\begin{equation}\label{eq:own-cube-exact-difference}
 V^-_{\rm restored}-V^-_{\rm old}
 =\sum_{P_4\text{ selected old first exit}}
          g_N(P_4)V_N^-(N^{\lambda_4}/P_4,p_4).
\end{equation}
Old contribution at this prefix was exactly zero: there is no old
residual lower density to subtract. Inserting a fictitious old
subtree and subtracting two one-sided estimates is unnecessary.

\paragraph{Actual support and the compiler.}
The first four full squares are checked in
\eqref{eq:own-cube-new-tests}. Residual lower well support gives
every later full square at level $\lambda_4$. For the selected
branch, old unselected slots are controlled by $G_I$. Every appended
prime is smaller than $p_4$; after all selected old slots have
occurred, a selected-remaining square is exactly its full square.
Thus the entire newly retained descendant vector satisfies the same
selected-product factorization gates, not merely the prefix product.

The signed lower compiler extends to this finite policy by its
existing proof, rather than by treating the old post-admission
definition as a theorem about every possible tree. The rough-projection
identity depends only on even first exits. The finite box rounding
uses inner tests on positive even coefficients and outer tests on
negative odd coefficients; these give coefficients below the raw
lower coefficients. Small-prime beta choices are lower for positive
and upper for negative coefficients. The added decision contributes
only finitely many affine walls at a fixed positive cutoff.
All restored atoms have at least four marked slots, so they use the
full/selected-square distribution interface just checked, not a new
unverified low-rank shortcut. The fixed-label group partitions,
factorials and moving-residue estimate are unchanged.
As usual choose the finite policy and all buffers first, then the
small-prime and local-scale errors, then take $N\to\infty$.

The ordinary residual density equality, uniformly on compact $s>2$,
turns \eqref{eq:own-cube-exact-difference} into the lower-start gain
\begin{equation}\label{eq:own-cube-gain}
 \Delta_{4,\mathrm{restore}}=
 4\int_{\mathcal B}
 \frac{\bar f((\lambda_4-S_4)/t_4)}
 {t_1t_2t_3t_4^2}\,dt_1dt_2dt_3dt_4,
 \qquad \bar f=f/(2e^\gamma).
\end{equation}
The factor four is the same normalization as in the old fourth
increment. An $L_2$ starting improvement carries Chen weight $1/4$.
The present box is above $a=.0834$ and is not an $L_1$ box.

\paragraph{The whole restored domain stays below the first delay break.}
There is a useful general simplification, not just for the box below.
Suppose $t_1\le h$, $t_1+3t_2<\nu$, $t_3<t_2$, and
$\nu\le C_3\le M\le\Lambda$. On the newly restored domain,
\[
 C_3\le S_3+3r<M,\qquad r=t_4,
\]
one has
\[
 S_3<\frac{2\nu+h}{3},\qquad
 r>\frac{\nu-h}{9}=\frac{22613}{504000},
\]
and therefore
\[
 0<s_{\rm new}-2
 =\frac{M-S_3-3r}{r}
 \le\frac{\Lambda-\nu}{r}
 <\frac{9(\Lambda-\nu)}{\nu-h}
 =\frac{84375}{90452}<1.
\]
Thus $2<s_{\rm new}<3$ throughout this domain: the exact elementary
formula $\bar f(s)=\log(s-1)/s$ applies, with no delay-table backend.
For a buffered baseline $\nu_\rho=\nu-\rho$, repeat the same proof
with its actual earlier history and $C_3\ge\nu_\rho$.
The strict comparison remains valid for
\[
 0<\rho<\frac{10\nu-9\Lambda-h}{10}
           =\frac{6077}{2240000},
\]
since $10\nu-9\Lambda-h=6077/224000>0$.

\paragraph{One exact rational restored box.}
Let $\mathcal B=\prod_{j=1}^4[l_j,u_j]$, with
\[
 (l_1,l_2,l_3,l_4)=\frac1{10000}(1399,1299,999,679),
 \quad
 (u_1,u_2,u_3,u_4)=\frac1{10000}(1401,1301,1001,681).
\]
It is strictly ordered and $u_1<h$. Its old second-cube margin is
\[
 \nu-u_1-3u_2=\frac{7113}{280000}>0.
\]
An upper bound for the actual third envelope is obtained by including
all seven nonempty selected-$v$ expressions, the three $W$ upper
planes $\Lambda-3t_j/8$, and inheritance $\nu$. Direct rational
range calculation shows their maximum on this box is
\[
 C_{3,+}=\frac{1+u_1}{2}=\frac{11401}{20000}.
\]
The leading single branch attains $C_3=(1+t_1)/2$ here; the
unconditional upper bound suffices for what follows. Every existing
buffered third cap lies below $C_{3,+}$. Meanwhile
\[
 S_{3,-}+3l_4=\frac{2867}{5000},\qquad
 S_{3,-}+3l_4-C_{3,+}=\frac{67}{20000}>0.
\]
Hence the entire box was rejected even by the ideal old $C_3$, and
also by any of its actual buffered ancestors. Those ancestor caps
are kept unchanged; there is no hidden cost of replacing them by a
smaller fixed cap.

Choose $I=\{2,4\}$ and $\lambda_4=593/1000$. On the whole box,
\[
 \inf v(t_2+t_4)=\frac{23733}{40000}>\lambda_4,
 \quad \max G_I=\frac{1077}{2000}<\lambda_4,
 \quad \max_{j\le4}(S_{j-1}+2t_j)=\frac{1013}{2000}<\lambda_4,
\]
and
\[
 \lambda_4-(S_{3,+}+3u_4)=\frac{23}{1250}>0.
\]
All branch, history and support inequalities have explicit fixed
slack. Its residual argument lies in
\[
 \frac{1546}{681}\le s\le\frac{222}{97},
\]
which is a compact subinterval of $(2,3)$. For $q\ge1$,
$\log q\ge2(q-1)/(q+1)$, because the derivative of the difference
is $(q-1)^2/[q(q+1)^2]\ge0$. Consequently
\[
 \bar f(s)\ge\frac{184/773}{222/97}=\frac{8924}{85803}.
\]
Using the exact volume $(1/5000)^4$ and the upper denominator
$u_1u_2u_3u_4^2$ in \eqref{eq:own-cube-gain} gives
\[
 \boxed{\displaystyle
 \Delta_{4,\mathrm{restore}}\ge
 \frac{5711360000}{72601602346142052183}
 >7.86\cdot10^{-11}.}
\]
The corresponding local Chen coefficient is at least
$1427840000/72601602346142052183$ before arbitrarily small finite
compilation errors. A strictly smaller positive value survives
those errors. This is a positivity witness, not a claim of a
material global improvement.

For comparison the originally proposed $I=\{1,4\}$ and
$\lambda_4=289/500$ are also licensed and give
$263840000/14807380504178267577$. These are alternatives on the
same domain: choose the better one, never add them.

\paragraph{Compatibility and present scope.}
The new box is disjoint from old $E_4$ and hence from the projection
of every previously counted old $E_5$ node. Preserve all old fourth
and fifth decisions on their old retained domain; restore only this
old-exit box and attach its ordinary lower residual, without a new
fifth upgrade. The raw difference is then exactly
\eqref{eq:own-cube-exact-difference}. There is no overlap with a
previous fourth or fifth gain. Any future continuous all-restored
policy needs the corresponding new admitted-domain approximation;
the old approximation statement with its fixed $E_4$ must not be
silently reinterpreted. This local finite box needs no new global
domain convergence theorem.

The companion standard-library-only Fraction script rechecks every
rational bound and the entire-domain $s<3$ inequality. It also
enumerates all $512$ divisibility patterns of a nine-symbol finite
tree, checks even first exits and both lower pointwise inequalities,
and exactly checks the raw-density difference. This finite toy is
a check on, not a substitute for, the general parity proof above.
No old file, numerical total, main source, research state,
formalization or publication is changed.
